```latex
\[
\newcommand{\R}{\mathbb{R}}
\newcommand{\diag}{\operatorname{diag}}
\]

For \(1\le r<\infty\), write
\[
\|A\|_{r\to r}:=\sup_{x\ne 0}\frac{\|Ax\|_r}{\|x\|_r}
\]
for the induced \(\ell^r\)-operator norm, and define
\[
\rho_r(A):=\inf\left\{\|DAD^{-1}\|_{r\to r}: D \text{ is positive diagonal}\right\}.
\]
Here positive diagonal means diagonal with strictly positive diagonal entries.

\[
\textbf{Part (a).}
\]

Let \(s=n-m\). We first prove the block norm identity
\[
\|\diag(A,0_s)\|_{r\to r}=\|A\|_{r\to r}
\]
for every \(A\in\R^{m\times m}\) and every \(s\ge 0\).

First note explicitly that
\[
\|A\|_{r\to r}
=
\sup_{\|u\|_r=1}\|Au\|_r.
\]
Indeed, if \(x\ne 0\), then \(u=x/\|x\|_r\) satisfies \(\|u\|_r=1\), and by homogeneity,
\[
\frac{\|Ax\|_r}{\|x\|_r}
=
\left\|A\frac{x}{\|x\|_r}\right\|_r
=
\|Au\|_r.
\]
Conversely, every \(u\) with \(\|u\|_r=1\) is an admissible nonzero vector in the original supremum. Since the \(\ell^r\)-unit sphere in \(\R^m\) is compact and \(u\mapsto \|Au\|_r\) is continuous, this supremum is actually attained.

Now take \(x=(x_1,x_2)\in \R^m\oplus \R^s\). Then
\[
\|x\|_r^r=\|x_1\|_r^r+\|x_2\|_r^r,
\]
so \(\|x_1\|_r\le \|x\|_r\). Hence
\[
\|\diag(A,0_s)x\|_r
=
\|(Ax_1,0)\|_r
=
\|Ax_1\|_r
\le
\|A\|_{r\to r}\|x_1\|_r
\le
\|A\|_{r\to r}\|x\|_r.
\]
Taking the supremum over \(x\ne 0\) gives
\[
\|\diag(A,0_s)\|_{r\to r}\le \|A\|_{r\to r}.
\]

For the reverse inequality, let \(u^\ast\in\R^m\) satisfy \(\|u^\ast\|_r=1\) and
\[
\|Au^\ast\|_r=\|A\|_{r\to r}.
\]
Then
\[
\|\diag(A,0_s)\|_{r\to r}
\ge
\frac{\|\diag(A,0_s)(u^\ast,0)\|_r}{\|(u^\ast,0)\|_r}
=
\frac{\|(Au^\ast,0)\|_r}{\|u^\ast\|_r}
=
\|Au^\ast\|_r
=
\|A\|_{r\to r}.
\]
Therefore
\[
\|\diag(A,0_s)\|_{r\to r}=\|A\|_{r\to r}.
\]

Now let \(D\in\R^{n\times n}\) be positive diagonal. With respect to the decomposition
\[
\R^n=\R^m\oplus \R^{n-m},
\]
we can write
\[
D=\diag(D_1,D_2),
\]
where \(D_1\in\R^{m\times m}\) and \(D_2\in\R^{(n-m)\times(n-m)}\) are positive diagonal. If \(n=m\), the \(D_2\)-block is absent.

Then
\[
D M_n D^{-1}
=
\diag(D_1,D_2)
\begin{pmatrix}
H&0\\
0&0_{n-m}
\end{pmatrix}
\diag(D_1^{-1},D_2^{-1})
=
\diag(D_1HD_1^{-1},0_{n-m}).
\]
By the block norm identity,
\[
\|D M_nD^{-1}\|_{r\to r}
=
\|D_1HD_1^{-1}\|_{r\to r}.
\]
As \(D\) ranges over all positive diagonal \(n\times n\) matrices, \(D_1\) ranges over all positive diagonal \(m\times m\) matrices, and the expression above is independent of \(D_2\). Hence
\[
\rho_r(M_n)
=
\inf_D \|D M_nD^{-1}\|_{r\to r}
=
\inf_{D_1}\|D_1HD_1^{-1}\|_{r\to r}
=
\rho_r(H).
\]
Thus
\[
\boxed{\rho_r(M_n)=\rho_r(H)}
\]
for every \(r\ge 1\).

\[
\textbf{Part (b).}
\]

We first prove the Hadamard formula needed below. Let \(H\in\R^{m\times m}\) be a Hadamard matrix, so every entry of \(H\) is \(\pm 1\) and
\[
HH^{\mathsf T}=mI_m.
\]
We claim that for every \(q\ge 2\),
\[
\rho_q(H)=m^{1-\frac1q}.
\]

First, from
\[
H\left(\frac1m H^{\mathsf T}\right)=I_m,
\]
we see that \(H\) has a right inverse. Since \(H\) is square, this implies \(H\) is invertible: indeed,
\[
1=\det(I_m)=\det(H)\det\left(\frac1mH^{\mathsf T}\right),
\]
so \(\det(H)\ne 0\). Therefore
\[
\frac1m H^{\mathsf T}=H^{-1}.
\]
Multiplying on the right by \(H\), we get
\[
\frac1m H^{\mathsf T}H=I_m,
\]
and hence
\[
H^{\mathsf T}H=mI_m.
\]

We shall use the standard norm comparison inequality: if \(1\le a\le q\), then for every \(z\in\R^m\),
\[
\|z\|_a\le m^{\frac1a-\frac1q}\|z\|_q.
\]
For \(a=q\) this is equality. For \(a<q\), Hölder's inequality gives
\[
\sum_{j=1}^m |z_j|^a
\le
\left(\sum_{j=1}^m |z_j|^q\right)^{a/q}
m^{1-a/q}.
\]
Taking \(a\)-th roots yields the desired inequality.

Let
\[
C:=m^{1-\frac1q}.
\]

For the upper bound, first suppose \(q=2\). Then for every \(x\in\R^m\),
\[
\|Hx\|_2^2
=
x^{\mathsf T}H^{\mathsf T}Hx
=
m\|x\|_2^2.
\]
Thus
\[
\|H\|_{2\to 2}=\sqrt m=m^{1-\frac12}.
\]

Now suppose \(q>2\). For \(x\in\R^m\), set \(y=Hx\). Since \(H^{\mathsf T}H=mI_m\),
\[
\|y\|_2=\|Hx\|_2=\sqrt m\,\|x\|_2.
\]
Using the norm comparison with \(a=2\),
\[
\|x\|_2\le m^{\frac12-\frac1q}\|x\|_q.
\]
Therefore
\[
\|y\|_2
\le
\sqrt m\,m^{\frac12-\frac1q}\|x\|_q
=
m^{1-\frac1q}\|x\|_q
=
C\|x\|_q.
\]

Also, since every entry of \(H\) is \(\pm1\),
\[
|(Hx)_i|
\le
\sum_{j=1}^m |x_j|
=
\|x\|_1.
\]
Hence
\[
\|y\|_\infty=\|Hx\|_\infty\le \|x\|_1.
\]
Using the norm comparison with \(a=1\),
\[
\|x\|_1\le m^{1-\frac1q}\|x\|_q=C\|x\|_q.
\]
Thus
\[
\|y\|_\infty\le C\|x\|_q.
\]

Now
\[
\|y\|_q^q
=
\sum_{i=1}^m |y_i|^q
=
\sum_{i=1}^m |y_i|^{q-2}|y_i|^2
\le
\|y\|_\infty^{q-2}\|y\|_2^2.
\]
Using the two bounds above,
\[
\|Hx\|_q^q
=
\|y\|_q^q
\le
(C\|x\|_q)^{q-2}(C\|x\|_q)^2
=
C^q\|x\|_q^q.
\]
Taking \(q\)-th roots gives
\[
\|Hx\|_q\le C\|x\|_q.
\]
Therefore
\[
\|H\|_{q\to q}\le m^{1-\frac1q}.
\]
Since \(I_m\) is an admissible positive diagonal matrix,
\[
\rho_q(H)\le \|H\|_{q\to q}\le m^{1-\frac1q}.
\]

For the reverse inequality, let
\[
D=\diag(d_1,\dots,d_m)
\]
be any positive diagonal matrix. Choose \(i\) such that
\[
d_i=\max_{1\le j\le m} d_j.
\]
Let
\[
u:=H^{\mathsf T}e_i.
\]
Because the entries of \(H\) are \(\pm1\), we have \(u\in\{\pm1\}^m\). Moreover,
\[
Hu=HH^{\mathsf T}e_i=me_i.
\]
Set
\[
x:=Du.
\]
Then \(x\ne 0\), and
\[
DHD^{-1}x
=
DHD^{-1}Du
=
DHu
=
D(me_i)
=
md_i e_i.
\]
Thus
\[
\|DHD^{-1}x\|_q=md_i.
\]
On the other hand,
\[
\|x\|_q
=
\|Du\|_q
=
\left(\sum_{j=1}^m d_j^q|u_j|^q\right)^{1/q}
=
\left(\sum_{j=1}^m d_j^q\right)^{1/q}
\le
\left(\sum_{j=1}^m d_i^q\right)^{1/q}
=
m^{1/q}d_i.
\]
Therefore
\[
\|DHD^{-1}\|_{q\to q}
\ge
\frac{\|DHD^{-1}x\|_q}{\|x\|_q}
\ge
\frac{md_i}{m^{1/q}d_i}
=
m^{1-\frac1q}.
\]
Since this holds for every positive diagonal \(D\),
\[
\rho_q(H)\ge m^{1-\frac1q}.
\]
Combining the upper and lower bounds,
\[
\boxed{\rho_q(H)=m^{1-\frac1q}}
\]
for every \(q\ge 2\).

It remains to construct the matrices needed for \(n>6\).

Define
\[
H_4=
\begin{pmatrix}
1&1&1&1\\
1&-1&1&-1\\
1&1&-1&-1\\
1&-1&-1&1
\end{pmatrix}.
\]
A direct computation gives
\[
H_4H_4^{\mathsf T}
=
4I_4,
\]
because each row has squared norm \(4\), and distinct rows have inner product \(0\). Hence \(H_4\) is a \(4\times 4\) Hadamard matrix.

Next define
\[
H_8:=
\begin{pmatrix}
H_4&H_4\\
H_4&-H_4
\end{pmatrix}.
\]
All entries of \(H_8\) are \(\pm1\), and
\[
H_8H_8^{\mathsf T}
=
\begin{pmatrix}
H_4&H_4\\
H_4&-H_4
\end{pmatrix}
\begin{pmatrix}
H_4^{\mathsf T}&H_4^{\mathsf T}\\
H_4^{\mathsf T}&-H_4^{\mathsf T}
\end{pmatrix}.
\]
Multiplying the block matrices,
\[
H_8H_8^{\mathsf T}
=
\begin{pmatrix}
2H_4H_4^{\mathsf T}&0\\
0&2H_4H_4^{\mathsf T}
\end{pmatrix}
=
\begin{pmatrix}
8I_4&0\\
0&8I_4
\end{pmatrix}
=
8I_8.
\]
Thus \(H_8\) is an \(8\times 8\) Hadamard matrix.

Now let \(n>6\) and let \(p>2\) be an integer. Then \(p\ge 3\), so both \(p\) and \(p+1\) are at least \(2\), and the Hadamard formula applies.

\[
\textbf{Case 1: } n=7.
\]
Take
\[
M=\diag(H_4,0_3)\in\R^{7\times 7}.
\]
By part \((a)\),
\[
\rho_p(M)=\rho_p(H_4),
\qquad
\rho_{p+1}(M)=\rho_{p+1}(H_4).
\]
Using the formula \(\rho_q(H_4)=4^{1-1/q}\), we obtain
\[
\rho_p(M)=4^{1-\frac1p},
\qquad
\rho_{p+1}(M)=4^{1-\frac1{p+1}}.
\]
Since
\[
\frac1p>\frac1{p+1},
\]
we have
\[
1-\frac1p<1-\frac1{p+1}.
\]
Because \(4>1\), the function \(t\mapsto 4^t\) is strictly increasing. Hence
\[
4^{1-\frac1p}<4^{1-\frac1{p+1}},
\]
and therefore
\[
\rho_p(M)<\rho_{p+1}(M).
\]

\[
\textbf{Case 2: } n\ge 8.
\]
Take
\[
M=\diag(H_8,0_{n-8})\in\R^{n\times n},
\]
with the zero block omitted when \(n=8\). By part \((a)\),
\[
\rho_p(M)=\rho_p(H_8),
\qquad
\rho_{p+1}(M)=\rho_{p+1}(H_8).
\]
Using the formula \(\rho_q(H_8)=8^{1-1/q}\), we get
\[
\rho_p(M)=8^{1-\frac1p},
\qquad
\rho_{p+1}(M)=8^{1-\frac1{p+1}}.
\]
Again,
\[
1-\frac1p<1-\frac1{p+1},
\]
and since \(8>1\),
\[
8^{1-\frac1p}<8^{1-\frac1{p+1}}.
\]
Therefore
\[
\rho_p(M)<\rho_{p+1}(M).
\]

Combining the two cases, for every integer \(n>6\) and every integer \(p>2\), there exists \(M\in\R^{n\times n}\) such that
\[
\boxed{\rho_p(M)<\rho_{p+1}(M)}.
\]
This proves Conjecture \(3\).
```